% Das einzige Beispiel, das absichtlich warnt. Deutsch beugt das Wort, das ein
% Bindeglied regiert -- des Vertrags, des Beschlusses --, also kann keine Regel
% die Wendung aus einem Betreff im Nominativ bilden. Ohne `ref label' sagt das
% Paket das einmal je Artefakt; mit `ref label' schweigt es und setzt die
% Wendung, die dort steht.
\IfDocumentMetadataTF{}{\DocumentMetadata{}}
\documentclass[12pt,german]{article}
\usepackage{babel}
\usepackage{regulatory}

\newcommand\translation[2]{#2}

\refdocument[ohne-]{example2-nl}{%
    path=./,%
    author=E. Nijenhuis,%
    subject=Beispielurkunde%
}

\refdocument[mit-]{example1-nl}{%
    path=./,%
    author=E. Nijenhuis,%
    subject=Musterurkunde,%
    ref label=der Musterurkunde%
}

\hypersetup{pdftitle={Beispiel Deutsch}}

\begin{document}
    \begin{center}
        \Huge Beispiel Deutsch
    \end{center}

    \article{Verweise}\label{art:verweise}
    \begin{paras}
        \item \label{lid:ohne} \Aref{ohne-lid:lorem}
        \item \label{lid:zweimal} \Aref{ohne-lid:lorem2}
        \item \label{lid:mit} \Aref{mit-sub:lorem2}
    \end{paras}
\end{document}
